%La parte central del trabajo refiere a lo que es producción propia o aporte del
%proyecto de grado, incluyendo las decisiones tomadas. Por ejemplo, puede incluir
%los requerimientos, el análisis y el diseño de la solución. Si el proyecto tiene una
%implementación, debe describirse en términos de decisiones tomadas en ese sentido.
%Los detalles de programación se dejan para los anexos.

En esta sección se describen los pasos que llevamos luego de la revisión de antecedentes y obtención del ranking de herramientas.

Como primer paso nos enfocamos en realizar una prueba piloto con la herramienta ClickUp, que luego derivó en ajustes tanto del documento de entrada, como la construcción de una solución esperada y su traducción al inglés. Finalmente, la decisión de evaluar solamente la generación de historias producidas por esta herramienta.

\section{Prueba piloto: Click Up}

La herramienta ClickUp \cite{clickup} se utilizó como prueba piloto y primer acercamiento al análisis, ya que se considera una herramienta popular utilizada en la gestión de proyectos. Además, posee una integración con inteligencia artificial llamada Brain AI \cite{brainai}.

ClickUp es una herramienta comercial; no es posible acceder a su código e implementación, pero sí fue posible hacer uso de la prueba gratuita.

ClickUp lanzó Brain AI en el año 2024 \cite{brainai}, y permite utilizar agentes dentro de la aplicación. Según la documentación oficial, Brain AI se apoya en modelos avanzados y además permite conectar modelos externos desde la propia interfaz (p. ej., ChatGPT o Claude). \cite{clickupai,clickupai2} %La integración está construida con ChatGPT-4o y ChatGPT-4.1 \cite{clickupai}. Asimismo, también es posible utilizar otros modelos directamente, como ChatGPT GPT-5, ChatGPT GPT-5 Mini, Claude Sonnet 4, Claude Opus 4, Gemini 2.5 Pro y Gemini 2.5 Flash \cite{clickupai2}.

Se consideró presente la capacidad de generación automática de historias de usuario, dadas las capacidades de los agentes de procesar los datos e información dentro de un entorno de trabajo y crear tarjetas o tareas automáticamente. Además, en el blog \cite{clickupai3} se ejemplifica cómo generar historias de usuario dentro de la plataforma. El agente permite refinar, resumir o mejorar la salida de la IA, lo que posibilita verificar las salidas, en línea con los aspectos de interés definidos.

Para corroborar las capacidades, cada integrante se registró con un usuario y probó las características de Brain AI hasta que se alcanzó la cuota de uso paga.

Se utilizó el documento técnico de prueba, tanto pasándolo como archivo adjunto dentro del chat de Brain AI en la plataforma, como subiéndolo al espacio de trabajo y solicitando que el asistente lo utilizara como fuente para la creación de la instrucción definida en la sección anterior.

A partir de esto, se confirmó la capacidad de cubrir los aspectos de interés definidos (elaboración de historias, criterios y verificación de conflictos) a partir del archivo de prueba inconsistente. La herramienta permite verificar o refinar las historias antes de confirmar su creación. Además, es posible solicitar que incluya vinculación con las fuentes entre el archivo de requisitos y la historia de usuario. Dentro de los tipos de entrada, se probó con PDF y también con texto directamente. La salida puede ser estructurada según la instrucción que se le indique a la IA; sin embargo, con la instrucción básica utilizada, las historias son devueltas como texto, y es posible confirmar la creación de la tarea. No es posible descargar en diferentes formatos, pero sí copiarla para crear un documento fuera de la herramienta.

Las pruebas fueron realizadas una vez por cada integrante del equipo, para identificar cómo evaluar la calidad de cada capacidad y definir la mejor metodología posible. En todos los casos, se generaron siete historias de usuario completas, correspondientes a los mismos requerimientos definidos, con un nivel adecuado de coherencia y un grado de interacción principalmente automático, con la posibilidad de refinar la historia en caso de ser necesario.

Analizando las historias de usuario desde la perspectiva INVEST, se verificó que la mayoría cumple con ser Independiente, Negociable y Valiosa, aunque no Estimable, Pequeña ni Testeable. Esto se debe a que el documento de prueba es superficial y no indaga en detalles que permitirían mejorar la historia. Por ejemplo, en la historia generada de Historia Clínica Electrónica: ``Como médico autorizado, quiero acceder y registrar diagnósticos, tratamientos y evolución en la historia clínica electrónica de mis pacientes para dar seguimiento adecuado. Criterios de aceptación:
Solo personal autorizado puede acceder a la historia clínica.
Se registra un log de accesos y modificaciones.
Se pueden añadir diagnósticos, tratamientos y estudios.
'' No se detallan la cantidad de información a registrar ni las reglas de validación de acceso, lo que dificulta estimar el esfuerzo requerido (Estimable) y diseñar pruebas objetivas (Testable). Además, la historia agrupa varias funcionalidades, como acceso, registro de diagnósticos, tratamientos, estudios y logging, por lo que no es pequeña, a menos que se subdivida en tareas. %(Sugerencia de respaldo: agregar cita breve a INVEST/Agile Alliance aquí para dar fundamento metodológico.)

En cuanto a los criterios de aceptación, evaluados bajo el enfoque SMART, se observó que generalmente no son muy Específicos, Medibles ni Alcanzables, pero sí Relevantes. Los criterios carecen de detalles precisos que permitan una validación objetiva. Por ejemplo, para la Gestión de Pacientes, el criterio ``Se puede crear un nuevo paciente ingresando nombre, documento, contacto y antecedentes médicos'' no especifica qué datos exactos son obligatorios, ni qué formato o validaciones deben cumplir (como longitud del documento, tipos de datos o validación de contacto). Sin embargo, este nivel de detalle tampoco estaba incluido en el documento de entrada, por lo que resulta difícil evaluar este aspecto de los criterios de aceptación con la fuente utilizada. %(Sugerencia de respaldo: agregar referencia a SMART original o guía académica que se use en el resto del informe.)

En cuanto a la detección de inconsistencias, se observó un desempeño aceptable, ya que el agente logró identificar correctamente los conflictos presentes en los documentos, aunque en algunos casos marcó como inconsistencias situaciones que solo correspondían a falta de detalle, evidenciando cierta sensibilidad excesiva en la interpretación. Por ejemplo, para la historia de Facturación y Cobro:
`` Como responsable de facturación, quiero generar facturas electrónicas y notificar a los pacientes para asegurar el cobro de los servicios. Criterios de aceptación:
Se puede generar una factura electrónica por cada servicio prestado.
El paciente recibe la notificación de la factura emitida.'' Detectó como inconsistencia:
``No se aclara el canal de notificación a pacientes (email, SMS, app).''.

Asimismo, se verificó que la herramienta puede mantener vinculación con las fuentes entre las historias generadas y el documento fuente, agregando referencias a las secciones de origen cuando se le solicita. Por ejemplo, al final de la historia agrega correctamente:``Referencia: Sección 3.1 del documento.'' También se comprobó que permite la creación de subtareas y la vinculación con épicas, aunque estas acciones deben ser indicadas explícitamente por el usuario, ya que no se realizan de manera automática.

%En general, los resultados confirman que ClickUp Brain AI posee capacidades sólidas para la elaboración de historias de usuario a partir de documentos de requisitos, con integración contextual, vinculación con las fuentes y verificación, aunque requiere ajustes en las instrucciones e interacción con el analista para mejorar la calidad y consistencia de las salidas generadas.



\section{Mejora de la entrada}\label{sec:mejora-entrada}

Dada la prueba piloto de ClickUp, se decidió agregar detalles faltantes al documento de requerimientos, agregando detalles como los Actores participantes del sistema.

Además, se revisó el prompt, agregando instrucciones como la plantilla de las historias de usuario, instruyendo el contenido que debe incluir para su verificación, evitando ambigüedades. Se mantuvo un enfoque de \textit{zero-shot prompting} (ver Sección~\ref{sec:prompt}), sin incluir ejemplos previos de historias de usuario, y se aplicó \textit{answer engineering} al definir una plantilla de salida y campos obligatorios que estructuran la respuesta del modelo.

\begin{verbatim}
Genera historias de usuario a partir del siguiente documento de requisitos.

Identificá los requerimientos funcionales del documento y redacta historias de usuario
en formato estándar:
“Como <rol>, quiero <objetivo>, para <beneficio>”.

Las historias deben ser claras y concisas.

Para cada historia incluí además:
- criterios de aceptación básicos
- indicar si la historia debería considerarse una épica (sí/no)
- referencia al fragmento del documento utilizado para generarla

Si detectas inconsistencias, ambigüedades o contradicciones en el documento,
indica explícitamente en qué parte se encuentran.

Adjunto el documento
\end{verbatim}

\section{Generación de historias esperadas}

Luego de la prueba piloto, surgieron dudas sobre cómo corroborar la completitud de la solución de manera formal. Si bien se tiene como marco de referencia el documento de requisitos y algunas de las funcionalidades principales que cubre, aún surgieron dudas de cómo corroborar que efectivamente todas las funcionalidades descritas en el documento eran cubiertas. Para esto se decidió analizar por nuestro lado y generar las historias de usuario nosotros, como analistas expertos y utilizar la solución como referencia.

Las historias fueron revisadas y pasaron por un proceso de reescritura al detectar detalles faltantes o posibles épicas que fueron divididas. De esta manera, construimos una solución esperada lo más completa posible para la comparación.

La versión final de las historias esperadas se encuentra en la Tabla \ref{tab:historias_esperadas} y sus criterios en la Tabla \ref{tab:criterios_esperados}.

\section{Traducción al Inglés}
El documento de requisitos, las historias esperadas y los prompts fueron redactados en español. Sin embargo, se decidió crear una versión en inglés, dado que los modelos reciben un mayor entrenamiento en inglés \cite{zhang2023multilingualLLM} y suelen tener mejores resultados. En primer lugar se tradujo el documento de requisitos, con ayuda de ChatGPT, definiendo así un vocabulario, para luego traducir las historias esperadas de manera consistente.

El prompt fue traducido al inglés y se muestra en la Figura \ref{fig:prompt_simple}.

\begin{figure}[h!]
\begin{verbatim}
Generate user stories from the following requirements document.
Identify the functional requirements in the document and write user stories
in standard format:
"As a <role>, I want to <goal>, so that <benefit>."
The stories should be clear and concise.
For each story, also include:
- basic acceptance criteria
- indicate if the story should be considered an epic (yes/no)
- reference to the section of the document used to generate it
If you detect inconsistencies, ambiguities, or contradictions in the document,
explicitly indicate where they are found.
I have attached the document.
\end{verbatim}
\caption{Prompt simple utilizado para la generación de historias de usuario}
\label{fig:prompt_simple}
\end{figure}

\section{Tipos de evaluación}
\label{central-enfoque}

Con el objetivo de formalizar la evaluación, se analizó el artículo \cite{yu2026evaluating_genai_quality_multicase}, donde se evalúan salidas generadas por inteligencia artificial para un caso de uso definido. En este artículo se nombran diferentes metodologías, separadas por manual, automática o asistida por inteligencia artificial, que pueden ser utilizadas para validar diferentes aspectos de la salida generada, como la completitud, la exactitud, la relevancia, etc. Esto se puede ver en la Tabla \ref{table:aritculo-eval} extraída del artículo.

\begin{longtable} {|p{2.5cm} | p{2.5cm} | p{2.5cm} | p{2.5cm}|}
\caption{Tabla 12 extraída de Evaluating the quality of GenAI applications in software
engineering: a multi-case study: Tipos de evaluación para las métricas identificadas}
\label{table:aritculo-eval} \\
\hline
\textbf{Métrica} & \textbf{Evaluación manual} & \textbf{Evaluación automatizada} & \textbf{Evaluación asistida por IA} \\
\hline

\textbf{Exactitud (Acc.)} & 
Sí (revisión manual) & 
Sí (verificaciones automáticas) & 
No \\

\textbf{Corrección (Corr.)} & 
Sí (verificación manual) & 
Sí (validación automatizada) & 
No \\

\textbf{Alineación (Align.)} & 
Sí (evaluación de expertos) & 
Sí (análisis semántico) & 
Sí (modelos de IA) \\

\textbf{Completitud (Comp.)} & 
Sí (verificación manual) & 
Sí (verificaciones automáticas) & 
No \\

\textbf{Relevancia (Rel.)} & 
Sí (revisión subjetiva) & 
No & 
Sí (puntuación con IA) \\

\textbf{Factualidad (Fac.)} & 
Sí (verificación de hechos) & 
No & 
Sí (verificación con IA) \\

\textbf{Seguridad (Sec.)} & 
Sí (revisión de seguridad) & 
Sí (escaneo de vulnerabilidades) & 
No \\

\textbf{Transparencia (Tran.)} & 
Sí (revisión de claridad) & 
No & 
No \\

\hline
\end{longtable}

Para el caso de uso de generación de documentos utilizando inteligencia artificial, el artículo define cómo se mide cada aspecto de calidad, ejemplificado en la Tabla \ref{table:aritculo-eval2} extraída también del artículo. Este caso de uso se aproxima a la generación automática de historias de usuario, ya que el texto es generado por la inteligencia artificial a partir de un documento.

\newpage
\begin{longtable}{|p{3cm}| p{4cm} |p{4cm}|}
\caption{Tabla 14 extraída de Evaluating the quality of GenAI applications in software
engineering: a multi-case study: Métricas y prácticas sugeridas para la generación de documentos}
\label{table:aritculo-eval2}\\
\hline
\textbf{Aspecto de calidad} & \textbf{Métrica} & \textbf{Pasos para aplicarla en la práctica} \\
\hline
\endfirsthead

\caption{Tabla 14 extraída de Evaluating the quality of GenAI applications in software
engineering: a multi-case study: Métricas y prácticas sugeridas para la generación de documentos - Continuación} \\
\hline
\textbf{Aspecto de calidad} & \textbf{Métrica} & \textbf{Pasos para aplicarla en la práctica} \\
\hline
\endhead

\textbf{Exactitud} &
Tasa de corrección específica de la tarea (basada en reglas) &
Paso 1: Definir el contenido requerido (por ejemplo, los tres cambios principales). \newline
Paso 2: Contar los elementos capturados frente a los omitidos o alucinados. \newline
Paso 3: Utilizar listas de verificación o plantillas para una revisión estructurada. \\

\hline

\textbf{Relevancia} &
BERTScore (similitud semántica basada en embeddings) &
Paso 1: Comparar la salida de la IA generativa con ejemplos citados. \newline
Paso 2: Establecer rangos de umbral según el tipo de tarea. \newline
Paso 3: Aplicar en modo batch para múltiples resúmenes. \\

\hline

\textbf{Factualidad} &
FActScore (coincidencia basada en embeddings y a nivel de entidades) &
Paso 1: Aplicar en contenido de alto riesgo (por ejemplo, documentos de cumplimiento). \newline
Paso 2: Automatizar la comparación con documentos fuente confiables si es posible. \newline
Paso 3: Agregar revisión manual para casos límite. \\

\hline

\textbf{Alineación con las expectativas} &
Tarjeta de evaluación del revisor (lista de verificación o puntuación humana) &
Paso 1: Solicitar a los interesados que evalúen claridad, cobertura y utilidad. \newline
Paso 2: Utilizar escalas tipo Likert o formularios de puntuación basados en tareas. \newline
Paso 3: Recopilar retroalimentación para mejorar los prompts y el flujo de revisión. \\

\hline
\end{longtable}

% Esto capaz que se puede mejorar con lo que se habia escrito antes, sobre porque se descartaron algunos de los algoritmos del articulo
Se analizaron las técnicas automáticas propuestas en el artículo y se decidió que podría ser posible utilizar BERTScore para evaluar la similitud semántica de las historias generadas con respecto a una solución esperada o el documento de requisitos. Otros algoritmos automáticos como FActScore fueron descartados por ser más estrictos y no aplicables al dominio de historias de usuario, donde es posible utilizar sinónimos o inferir requerimientos de los documentos de requisitos. Teniendo esto en cuenta, se decidió definir un proceso completo de evaluación para historias de usuario, especificando evaluación manual, automática o asistida, definidos en el Capítulo \ref{cap:evaluacion}.